\cleardoublepage{}
\begin{center}
    \bfseries \zihao{3} 致~谢
\end{center}

时光荏苒，数载求学生涯即将画上句点。
当我坐下来写这篇致谢时，脑海中涌现出无数个曾经的瞬间——
深夜实验室亮着的灯、反复调试却迟迟收敛的训练曲线、
以及那些在困惑与突破之间交替度过的漫长午后。
这一路走来，我从未独行，每一步都有人给予光亮。

首先，诚挚感谢我的导师金小刚教授。
从论文选题到论文完成，金老师始终以开阔的学术视野为我指引方向。
金老师对研究问题有着敏锐的直觉与深刻的洞察，
每当我在方向选择上感到迷茫，一次谈话往往便能拨云见日。
能够在金老师的培养下完成这段学习历程，我深感荣幸与感激。

特别感谢带我走进科研世界的张宇晴学姐。
刚入组时，我对科研几乎一无所知，
是学姐手把手地教我阅读文献、修改代码、设计实验、分析结果。
学姐从不吝惜时间，无论是深夜的一条消息还是突如其来的疑问，
总能得到她耐心细致的回复。
在方法实现和实验设计的细节上，学姐亦不厌其烦地悉心指导。

感谢浙江大学多年来为我提供的广阔平台与丰厚资源。
这里有严谨求是的学术传统，有来自不同领域的优秀师友，
也有无数次让我眼界大开的讲座与交流。
浙大给予我的，不仅是专业知识的积累，
更是独立人格与批判思维的塑造。
能够在这里度过最重要的求学岁月，是我一生的财富。

最后，我要将最深的感谢献给我的父母。
这些年，无论我遇到怎样的困难与挫折，
他们从未施加任何压力，只是默默地陪伴与支持。
同样感谢我的女朋友陈子逸。
科研的道路并非总是一帆风顺，
每当我陷入焦虑与自我怀疑，
她总能在第一时间察觉到我的情绪，
用温柔而坚定的话语将我从低谷中拉回。
日常生活里她的陪伴与照料，
让紧绷的科研生活多了许多温度与色彩。
有她在身边，是我这段时光里最幸运的事之一。

谨以此文，献给所有在这段旅程中给予我帮助与温暖的人。